% Manual de usuario del proyecto final de Computacion Grafica
% Se compila con latexmk (ver latexmkrc). Las capturas se toman de ../docs/evidencias/ y de img/
% Lo que falta se marca con \pendiente{...}
\documentclass[11pt,a4paper]{article}
% Estilo comun de los manuales, el mismo de la propuesta y del reporte de avances:
% fondo oscuro, titulos morados y recuadro de encabezado. Se carga con % Estilo comun de los manuales, el mismo de la propuesta y del reporte de avances:
% fondo oscuro, titulos morados y recuadro de encabezado. Se carga con % Estilo comun de los manuales, el mismo de la propuesta y del reporte de avances:
% fondo oscuro, titulos morados y recuadro de encabezado. Se carga con \input{estilo}.
\usepackage[utf8]{inputenc}
\usepackage[T1]{fontenc}
\usepackage[spanish,es-noshorthands,es-tabla]{babel}
\usepackage{lmodern}
\usepackage[margin=2.5cm]{geometry}
\usepackage{graphicx}
\usepackage{subcaption}
\usepackage{float}
\usepackage[table]{xcolor}
\usepackage{titlesec}
\usepackage{tcolorbox}
\usepackage{enumitem}
\usepackage{array}
\usepackage{longtable}
\usepackage{amssymb}
\usepackage{pgfgantt}
\usepackage{hyperref}

\definecolor{fondo}{HTML}{1E1E2E}
\definecolor{texto}{HTML}{CDCDD6}
\definecolor{morado}{HTML}{B48EDB}
\definecolor{azul}{HTML}{8AADF4}
\definecolor{rosa}{HTML}{E3A8D6}
\definecolor{verde}{HTML}{A6DA95}
\definecolor{linea}{HTML}{55556A}

\pagecolor{fondo}
\color{texto}
\hypersetup{colorlinks=true, linkcolor=azul, urlcolor=azul, citecolor=azul}
\arrayrulecolor{texto}
\graphicspath{{../docs/evidencias/}{img/}}

\setlength{\parindent}{0pt}
\setlength{\parskip}{6pt}

% secciones numeradas, moradas y con una linea abajo
\titleformat{\section}{\Large\bfseries\color{morado}}{\thesection.}{8pt}{}[{\color{linea}\titlerule}]
\titleformat{\subsection}{\large\bfseries\color{azul}}{\thesubsection}{8pt}{}
\titleformat{\subsubsection}{\normalsize\bfseries\color{rosa}}{}{0pt}{}
\titlespacing*{\section}{0pt}{18pt}{10pt}
\titlespacing*{\subsection}{0pt}{12pt}{6pt}

\setlist[itemize,1]{label=\scriptsize$\blacksquare$}
\setlist{itemsep=3pt}

% lo que falta por escribir o por confirmar; se busca con "pendiente" para no dejar ninguno
\newcommand{\pendiente}[1]{%
  \begin{tcolorbox}[colback=fondo, colframe=rosa, coltext=rosa, arc=2mm, boxrule=0.8pt, left=6pt, right=6pt, top=4pt, bottom=4pt]
    \textbf{Pendiente:} #1
  \end{tcolorbox}}

% tecla con recuadro
\newcommand{\tecla}[1]{\tcbox[on line, colback=fondo, colframe=azul, coltext=azul, arc=1mm, boxrule=0.6pt, left=2pt, right=2pt, top=0pt, bottom=0pt]{\small\ttfamily #1}}

\newcommand{\dosfotos}[5]{%
  \begin{figure}[H]
    \centering
    \begin{subfigure}{0.48\textwidth}
      \includegraphics[width=\linewidth]{#1}
      \caption{#2}
    \end{subfigure}\hfill
    \begin{subfigure}{0.48\textwidth}
      \includegraphics[width=\linewidth]{#3}
      \caption{#4}
    \end{subfigure}
    \caption{#5}
  \end{figure}}

\newcommand{\unafoto}[2]{%
  \begin{figure}[H]
    \centering
    \includegraphics[width=0.8\textwidth]{#1}
    \caption{#2}
  \end{figure}}

% encabezado en recuadro: titulo del documento
\newcommand{\encabezado}[1]{%
  \begin{tcolorbox}[colback=fondo, colframe=morado, coltext=texto, arc=4mm, boxrule=0.8pt, left=10pt, right=10pt, top=8pt, bottom=8pt]
    \centering
    {\Large\bfseries\color{azul} #1 $\|$ Calenda de Oaxaca}\\[4pt]
    {\large\bfseries Facultad de Ingeniería -- UNAM}\\[10pt]
    \begin{tabular}{@{}p{0.47\linewidth}p{0.47\linewidth}@{}}
      {\color{rosa}\textbf{Equipo:}} 16 & {\color{rosa}\textbf{Materia:}} Computación Gráfica e I.H.C. \\
      {\color{rosa}\textbf{Número de cuenta:}} 421112601 & {\color{rosa}\textbf{Grupo:}} 01 \\
      {\color{rosa}\textbf{Semestre:}} 2027-1 & {\color{rosa}\textbf{Entrega:}} 18 de noviembre de 2026 \\
    \end{tabular}
  \end{tcolorbox}}
.
\usepackage[utf8]{inputenc}
\usepackage[T1]{fontenc}
\usepackage[spanish,es-noshorthands,es-tabla]{babel}
\usepackage{lmodern}
\usepackage[margin=2.5cm]{geometry}
\usepackage{graphicx}
\usepackage{subcaption}
\usepackage{float}
\usepackage[table]{xcolor}
\usepackage{titlesec}
\usepackage{tcolorbox}
\usepackage{enumitem}
\usepackage{array}
\usepackage{longtable}
\usepackage{amssymb}
\usepackage{pgfgantt}
\usepackage{hyperref}

\definecolor{fondo}{HTML}{1E1E2E}
\definecolor{texto}{HTML}{CDCDD6}
\definecolor{morado}{HTML}{B48EDB}
\definecolor{azul}{HTML}{8AADF4}
\definecolor{rosa}{HTML}{E3A8D6}
\definecolor{verde}{HTML}{A6DA95}
\definecolor{linea}{HTML}{55556A}

\pagecolor{fondo}
\color{texto}
\hypersetup{colorlinks=true, linkcolor=azul, urlcolor=azul, citecolor=azul}
\arrayrulecolor{texto}
\graphicspath{{../docs/evidencias/}{img/}}

\setlength{\parindent}{0pt}
\setlength{\parskip}{6pt}

% secciones numeradas, moradas y con una linea abajo
\titleformat{\section}{\Large\bfseries\color{morado}}{\thesection.}{8pt}{}[{\color{linea}\titlerule}]
\titleformat{\subsection}{\large\bfseries\color{azul}}{\thesubsection}{8pt}{}
\titleformat{\subsubsection}{\normalsize\bfseries\color{rosa}}{}{0pt}{}
\titlespacing*{\section}{0pt}{18pt}{10pt}
\titlespacing*{\subsection}{0pt}{12pt}{6pt}

\setlist[itemize,1]{label=\scriptsize$\blacksquare$}
\setlist{itemsep=3pt}

% lo que falta por escribir o por confirmar; se busca con "pendiente" para no dejar ninguno
\newcommand{\pendiente}[1]{%
  \begin{tcolorbox}[colback=fondo, colframe=rosa, coltext=rosa, arc=2mm, boxrule=0.8pt, left=6pt, right=6pt, top=4pt, bottom=4pt]
    \textbf{Pendiente:} #1
  \end{tcolorbox}}

% tecla con recuadro
\newcommand{\tecla}[1]{\tcbox[on line, colback=fondo, colframe=azul, coltext=azul, arc=1mm, boxrule=0.6pt, left=2pt, right=2pt, top=0pt, bottom=0pt]{\small\ttfamily #1}}

\newcommand{\dosfotos}[5]{%
  \begin{figure}[H]
    \centering
    \begin{subfigure}{0.48\textwidth}
      \includegraphics[width=\linewidth]{#1}
      \caption{#2}
    \end{subfigure}\hfill
    \begin{subfigure}{0.48\textwidth}
      \includegraphics[width=\linewidth]{#3}
      \caption{#4}
    \end{subfigure}
    \caption{#5}
  \end{figure}}

\newcommand{\unafoto}[2]{%
  \begin{figure}[H]
    \centering
    \includegraphics[width=0.8\textwidth]{#1}
    \caption{#2}
  \end{figure}}

% encabezado en recuadro: titulo del documento
\newcommand{\encabezado}[1]{%
  \begin{tcolorbox}[colback=fondo, colframe=morado, coltext=texto, arc=4mm, boxrule=0.8pt, left=10pt, right=10pt, top=8pt, bottom=8pt]
    \centering
    {\Large\bfseries\color{azul} #1 $\|$ Calenda de Oaxaca}\\[4pt]
    {\large\bfseries Facultad de Ingeniería -- UNAM}\\[10pt]
    \begin{tabular}{@{}p{0.47\linewidth}p{0.47\linewidth}@{}}
      {\color{rosa}\textbf{Equipo:}} 16 & {\color{rosa}\textbf{Materia:}} Computación Gráfica e I.H.C. \\
      {\color{rosa}\textbf{Número de cuenta:}} 421112601 & {\color{rosa}\textbf{Grupo:}} 01 \\
      {\color{rosa}\textbf{Semestre:}} 2027-1 & {\color{rosa}\textbf{Entrega:}} 18 de noviembre de 2026 \\
    \end{tabular}
  \end{tcolorbox}}
.
\usepackage[utf8]{inputenc}
\usepackage[T1]{fontenc}
\usepackage[spanish,es-noshorthands,es-tabla]{babel}
\usepackage{lmodern}
\usepackage[margin=2.5cm]{geometry}
\usepackage{graphicx}
\usepackage{subcaption}
\usepackage{float}
\usepackage[table]{xcolor}
\usepackage{titlesec}
\usepackage{tcolorbox}
\usepackage{enumitem}
\usepackage{array}
\usepackage{longtable}
\usepackage{amssymb}
\usepackage{pgfgantt}
\usepackage{hyperref}

\definecolor{fondo}{HTML}{1E1E2E}
\definecolor{texto}{HTML}{CDCDD6}
\definecolor{morado}{HTML}{B48EDB}
\definecolor{azul}{HTML}{8AADF4}
\definecolor{rosa}{HTML}{E3A8D6}
\definecolor{verde}{HTML}{A6DA95}
\definecolor{linea}{HTML}{55556A}

\pagecolor{fondo}
\color{texto}
\hypersetup{colorlinks=true, linkcolor=azul, urlcolor=azul, citecolor=azul}
\arrayrulecolor{texto}
\graphicspath{{../docs/evidencias/}{img/}}

\setlength{\parindent}{0pt}
\setlength{\parskip}{6pt}

% secciones numeradas, moradas y con una linea abajo
\titleformat{\section}{\Large\bfseries\color{morado}}{\thesection.}{8pt}{}[{\color{linea}\titlerule}]
\titleformat{\subsection}{\large\bfseries\color{azul}}{\thesubsection}{8pt}{}
\titleformat{\subsubsection}{\normalsize\bfseries\color{rosa}}{}{0pt}{}
\titlespacing*{\section}{0pt}{18pt}{10pt}
\titlespacing*{\subsection}{0pt}{12pt}{6pt}

\setlist[itemize,1]{label=\scriptsize$\blacksquare$}
\setlist{itemsep=3pt}

% lo que falta por escribir o por confirmar; se busca con "pendiente" para no dejar ninguno
\newcommand{\pendiente}[1]{%
  \begin{tcolorbox}[colback=fondo, colframe=rosa, coltext=rosa, arc=2mm, boxrule=0.8pt, left=6pt, right=6pt, top=4pt, bottom=4pt]
    \textbf{Pendiente:} #1
  \end{tcolorbox}}

% tecla con recuadro
\newcommand{\tecla}[1]{\tcbox[on line, colback=fondo, colframe=azul, coltext=azul, arc=1mm, boxrule=0.6pt, left=2pt, right=2pt, top=0pt, bottom=0pt]{\small\ttfamily #1}}

\newcommand{\dosfotos}[5]{%
  \begin{figure}[H]
    \centering
    \begin{subfigure}{0.48\textwidth}
      \includegraphics[width=\linewidth]{#1}
      \caption{#2}
    \end{subfigure}\hfill
    \begin{subfigure}{0.48\textwidth}
      \includegraphics[width=\linewidth]{#3}
      \caption{#4}
    \end{subfigure}
    \caption{#5}
  \end{figure}}

\newcommand{\unafoto}[2]{%
  \begin{figure}[H]
    \centering
    \includegraphics[width=0.8\textwidth]{#1}
    \caption{#2}
  \end{figure}}

% encabezado en recuadro: titulo del documento
\newcommand{\encabezado}[1]{%
  \begin{tcolorbox}[colback=fondo, colframe=morado, coltext=texto, arc=4mm, boxrule=0.8pt, left=10pt, right=10pt, top=8pt, bottom=8pt]
    \centering
    {\Large\bfseries\color{azul} #1 $\|$ Calenda de Oaxaca}\\[4pt]
    {\large\bfseries Facultad de Ingeniería -- UNAM}\\[10pt]
    \begin{tabular}{@{}p{0.47\linewidth}p{0.47\linewidth}@{}}
      {\color{rosa}\textbf{Equipo:}} 16 & {\color{rosa}\textbf{Materia:}} Computación Gráfica e I.H.C. \\
      {\color{rosa}\textbf{Número de cuenta:}} 421112601 & {\color{rosa}\textbf{Grupo:}} 01 \\
      {\color{rosa}\textbf{Semestre:}} 2027-1 & {\color{rosa}\textbf{Entrega:}} 18 de noviembre de 2026 \\
    \end{tabular}
  \end{tcolorbox}}


\begin{document}

\encabezado{Manual de usuario}

\tableofcontents
\clearpage

\section{Qué es este programa}

Es un recorrido en 3D por el centro de Oaxaca durante una calenda, el desfile con música, chinas oaxaqueñas, marmotas y toritos que recorre las calles para anunciar la fiesta de la Guelaguetza.
El recorrido empieza en el atrio del templo de Santo Domingo, baja por el andador Macedonio Alcalá y termina en el Zócalo, frente al kiosco.

Dentro del programa se puede caminar libremente por las calles, ver la escena desde cinco cámaras, cambiar entre la tarde y la noche, ver pasar el desfile y acercarse a los edificios y personajes para leer qué son y qué significan.

\unafoto{retoque_texturas/02_santo_domingo}{El templo de Santo Domingo de Guzmán desde la cámara de la fachada.}

\section{Qué se necesita}

\begin{itemize}
  \item Windows 10 u 11.
  \item Una tarjeta de video con OpenGL 4.3 o más reciente.
  \item Visual Studio 2026 (o 2022) con la carga de trabajo \textbf{Desarrollo para el escritorio con C++}.
\end{itemize}

\clearpage
\section{Cómo abrir y ejecutar el proyecto}

\subsection{Con Visual Studio}

\begin{enumerate}
  \item Descargar el proyecto de \url{https://github.com/chavez-inst/ProyectoFinalCompuGraficaTeoria} (botón \textbf{Code} y luego \textbf{Download ZIP}) y descomprimirlo en una carpeta.
  \item Descomprimir \texttt{Librerias.zip} dentro de la carpeta del proyecto, de forma que quede la carpeta \texttt{Librerias} junto a \texttt{src}, \texttt{Models} y \texttt{Textures}.
  \item Abrir \texttt{Proyecto\_Final.slnx}. Con Visual Studio 2022 se abre directamente \nolinkurl{Proyecto_Final.vcxproj}. Si Visual Studio pregunta si se quiere cambiar la versión de las herramientas de compilación, aceptar.
  \item Arriba, elegir la configuración \textbf{Debug} o \textbf{Release} y la plataforma \textbf{x64}.
  \item Presionar \tecla{F5} (o el botón verde de \textbf{Depurador local de Windows}).
\end{enumerate}

\subsection{Al abrir}

El programa abre en una ventana del tamaño del monitor. La escena empieza de tarde, con la cámara libre en el andador, frente al atrio de Santo Domingo y viendo hacia el sur.

\unafoto{final/01_vista_inicial}{La cámara libre al abrir el programa, sobre el andador y viendo hacia el sur.}

\clearpage
\section{Controles}

\begin{table}[H]
  \centering
  \begin{tabular}{|>{\centering\arraybackslash}p{0.2\linewidth}|p{0.7\linewidth}|}
    \hline
    \textbf{Tecla} & \textbf{Qué hace} \\ \hline
    \tecla{W} \tecla{A} \tecla{S} \tecla{D} & Caminar con la cámara libre: adelante, izquierda, atrás y derecha \\ \hline
    Mouse & Voltear la cámara libre hacia los lados y hacia arriba o abajo \\ \hline
    \tecla{Shift} & Mantener presionada junto con WASD para ir tres veces más rápido \\ \hline
    \tecla{1} a \tecla{5} & Ir directo a una cámara (ver la sección \ref{sec:camaras}) \\ \hline
    \tecla{C} & Pasar a la siguiente cámara \\ \hline
    \tecla{F} & Soltar una cámara fija: la cámara libre aparece en ese mismo lugar y desde ahí se puede caminar \\ \hline
    \tecla{N} & Cambiar entre tarde y noche \\ \hline
    \tecla{P} & Detener el desfile o hacer que siga caminando \\ \hline
    \tecla{R} & Regresar el desfile al atrio de Santo Domingo para verlo desde el principio \\ \hline
    \tecla{I} & Junto a una carta de lotería: tomarla, voltearla para leer su ficha y guardarla (ver la sección \ref{sec:cartas}) \\ \hline
    \tecla{L} & Mostrar u ocultar el tablero con las cartas que ya se encontraron \\ \hline
    \tecla{M} & Quitar o regresar todo el sonido \\ \hline
    \tecla{O} & Poner o quitar el modo Día de Muertos (ver la sección \ref{sec:muertos}) \\ \hline
    \tecla{ESC} & Cerrar el programa \\ \hline
  \end{tabular}
  \caption{Controles del recorrido.}
\end{table}

\section{Cámaras}
\label{sec:camaras}

\begin{table}[H]
  \centering
  \begin{tabular}{|>{\centering\arraybackslash}p{0.08\linewidth}|p{0.25\linewidth}|p{0.55\linewidth}|}
    \hline
    \textbf{Tecla} & \textbf{Cámara} & \textbf{Qué muestra} \\ \hline
    \tecla{1} & Libre & Se mueve con WASD y el mouse. No deja salir de la ciudad ni subir más de 30 m \\ \hline
    \tecla{2} & Fachada & La fachada de Santo Domingo desde el atrio, donde empieza la calenda \\ \hline
    \tecla{3} & Desfile & Va a un lado del desfile y lo acompaña del atrio al Zócalo \\ \hline
    \tecla{4} & Zócalo & La plaza en diagonal desde una esquina, con el kiosco al centro \\ \hline
    \tecla{5} & Aérea & Todo el recorrido desde 300 m de altura, con el norte hacia arriba \\ \hline
  \end{tabular}
  \caption{Las cinco cámaras.}
\end{table}

\dosfotos{final/12_camara_desfile}{Cámara del desfile (3)}{final/09_zocalo_tarde}{Cámara del Zócalo (4)}{Dos de las cámaras fijas.}

La vista aérea sirve como mapa: el atrio de Santo Domingo queda arriba a la derecha, el andador baja por el centro y el Zócalo está abajo.
Desde cualquier cámara fija se puede presionar \tecla{F} para seguir caminando desde ahí. Si se suelta la aérea, la cámara baja a 30 m, que es lo más alto que llega la libre.

\unafoto{retoque_texturas/05_aerea}{Cámara aérea (5): el recorrido completo, del atrio al Zócalo.}

\section{Tarde y noche}

Con \tecla{N} se cambia la hora. De tarde el sol está bajo y la cantera se ve naranja; de noche hay estrellas, la luz es azul y se encienden las farolas del andador y del Zócalo y los reflectores de los templos y del kiosco.

\dosfotos{final/09_zocalo_tarde}{Zócalo de tarde}{final/10_zocalo_noche}{Zócalo de noche}{El mismo lugar en los dos horarios.}

\section{Modo Día de Muertos}
\label{sec:muertos}

Con \tecla{O} la escena se convierte en una noche de Día de Muertos, y con otra \tecla{O} regresa a como estaba. El cielo se pone más oscuro y morado, las farolas bajan y casi toda la luz sale de las velas y de unos focos de color: Santo Domingo en morado, la Sangre de Cristo en turquesa y el kiosco en naranja. La luz de las velas parpadea.

En el andador cuelgan cortinas de cempasúchil, hay velas flotando entre el papel picado y este cambia por uno con calaveras. El kiosco del Zócalo se adorna con esferas y guirnaldas de cempasúchil, y a su lado hay un altar de varios escalones con un arco de flores, macetas, calaveritas, velas y un tapete de arena. También hay una catrina, perros xoloitzcuintles y tres ofrendas con pan, veladoras y calaveritas. Los espectadores y los del desfile llevan la cara pintada de calavera. Todo lo demás (el desfile, las cámaras, las cartas y el sonido) funciona igual.

\dosfotos{muertos/02_velas_flotantes}{Velas flotantes en el andador}{muertos/11_altar}{Altar junto al kiosco}{El modo Día de Muertos.}

\section{Sonido}

El programa tiene sonido en 3D: lo que se oye depende de dónde está la cámara activa. La música de la banda (\textit{Son Calenda} y \textit{Dios nunca muere}) sale del desfile, así que se oye fuerte cerca de él y más bajo al alejarse. De fondo se oye a la gente y, junto al torito, sus cohetes. Cuando el desfile sale del atrio suenan las campanas de Santo Domingo y truena un cohete; al llegar frente al kiosco truena otro.

Con \tecla{P} también se pausan la música y los cohetes del torito, con \tecla{R} la banda vuelve a la primera canción y con \tecla{M} se quita o se regresa todo el sonido.

\section{El recorrido}

Estos son los lugares principales, en el orden en que los recorre la calenda:

\begin{enumerate}
  \item \textbf{Atrio de Santo Domingo.} La explanada frente al templo, con jardineras de agaves, palmeras, árboles de flores rojas y dos puestos de artesanías: uno de barro y pan, y otro de alebrijes.
  \item \textbf{Andador Macedonio Alcalá.} La calle peatonal que une Santo Domingo con el Zócalo. Tiene casas de colores con macetas y bugambilias, y papel picado de un lado a otro de la calle.
  \item \textbf{Templo de la Sangre de Cristo.} Está a un lado del andador, con su fachada de cantera hacia la calle.
  \item \textbf{Jardín Labastida.} Una plaza pequeña con jacarandas, bugambilias y un puesto de textiles.
  \item \textbf{Zócalo.} La plaza principal, con el kiosco al centro, la fuente, jardineras, bancas de hierro, carritos de elotes y los portales alrededor.
\end{enumerate}

\section{El desfile}

En cuanto se abre el programa, la calenda sale de la puerta de Santo Domingo, cruza el atrio, da vuelta en el andador y baja por él hasta el Zócalo. Tarda unos dos minutos en llegar; frente al kiosco se detiene unos segundos y vuelve a empezar desde el atrio.

La mejor forma de verla es con la cámara \tecla{3}, que la acompaña todo el camino. Con \tecla{P} el desfile se detiene (los personajes se quedan quietos donde van) y con otra \tecla{P} sigue. Con \tecla{R} regresa al atrio para verlo desde el principio.

\unafoto{final/02_salida_atrio}{La calenda saliendo del atrio de Santo Domingo.}

La fila va en este orden:

\begin{enumerate}
  \item \textbf{Danzante con la marmota.} Abre el desfile. Lleva con las dos manos un palo largo con la marmota, un globo grande de tela con banderitas de colores, que gira y se mece de un lado a otro.
  \item \textbf{Cuatro chinas oaxaqueñas}, en dos filas, cada una con la falda de otro color. Llevan sobre la cabeza una canasta adornada con flores y la sostienen con la mano izquierda; con la derecha sacan un dulce y lo lanzan hacia la gente.
  \item \textbf{Tres músicos} con trompeta, que la levantan y la bajan mientras tocan.
  \item \textbf{El torito}, al final. Un danzante carga sobre la cabeza el armazón del toro, que se mece de lado a lado; a los costados tiene dos ruedas de cohetes que giran y echan chispas.
\end{enumerate}

\dosfotos{final/04_marmota}{Danzante con la marmota}{final/03_chinas_musicos}{Chinas oaxaqueñas y músicos}{Los personajes del desfile.}

\dosfotos{img/musicos_torito}{Músicos y, detrás, el torito}{img/espectadores_globos}{Espectador con globos}{El final de la fila y la gente que lo ve pasar.}

A los lados del andador, en el atrio y en el Zócalo hay espectadores que saludan con la mano y voltean la cabeza. Algunos llevan tres globos de colores que suben, bajan y se mecen.

\section{Información cultural: la lotería de la calenda}
\label{sec:cartas}

La información del recorrido está en diez cartas de lotería repartidas por la escena. Cada carta habla de un personaje o de un lugar y flota cerca de él, girando para que se note. De lejos se ve su reverso, con un signo de interrogación y una pista de dónde está.

\begin{itemize}
  \item \textbf{Cartas fijas}: la calenda en el atrio, Santo Domingo frente a su puerta, la Sangre de Cristo en la banqueta del andador, los alebrijes sobre su petate en el jardín Labastida, los tapetes junto a los puestos de textiles y el Zócalo frente al kiosco.
  \item \textbf{Cartas que van con el desfile}: la marmota, la primera china oaxaqueña, el músico de en medio y el torito. Se mueven con su personaje, así que hay que alcanzarlo (la cámara \tecla{3} ayuda) o detener el desfile con \tecla{P}.
\end{itemize}

\subsection{Cómo se toma una carta}

\begin{enumerate}
  \item Acercarse a la carta con cualquier cámara. Al estar cerca aparece abajo el aviso \textbf{I: tomar la carta}.
  \item Presionar \tecla{I}. La carta aparece en grande en el centro de la pantalla, con la figura de la lotería, y desaparece de la escena.
  \item Presionar \tecla{I} otra vez. La carta se voltea y muestra su ficha, con un texto corto sobre el tema. Al mismo tiempo se oye la narración de ese texto y la música baja para que se escuche.
  \item Presionar \tecla{I} una vez más para guardar la carta y seguir con el recorrido.
\end{enumerate}

\dosfotos{final/05_carta_marmota}{Carta de la marmota junto a su danzante, con el aviso para tomarla}{final/06_carta_frente}{Frente de la carta al tomarla}{Al acercarse a una carta y tomarla con la tecla I.}

\dosfotos{final/07_carta_ficha}{Ficha de la marmota, con el aviso para guardarla}{final/08_tablero}{Tablero con la primera carta encontrada}{La ficha de una carta y el tablero (tecla L).}

\subsection{El tablero}

Con \tecla{L} se abre el tablero, con lugar para las diez cartas. Las que ya se tomaron se ven de frente; las que faltan muestran su pista (por ejemplo, ``Echa chispas al final del desfile'' o ``Entre los puestos de textiles''). Otra \tecla{L} lo cierra. Cuando se juntan las diez cartas aparece el letrero de lotería y truena un cohete.

\end{document}
